\section{El juego continuo de trayectorias y la ejecución mientras se calcula}
\label{sec:continuacion}

El algoritmo anterior negocia rutas de un catálogo finito. El banco de
navegación con ciento veinte mundos aleatorizados \cite{mayorga2026compendio}
deja el diagnóstico de esa elección: decidir solo el \emph{dónde} (ruta) deja un
residuo de reparación del $58{,}3\%$; decidir \emph{dónde y cuándo}
(ruta--tiempo) lo baja al $26{,}7\%$; pero revisar dentro del mismo catálogo con
orden $h\le3$ resuelve solo $1$ de los $32$ residuos restantes. La limitación ya
no es la profundidad de la revisión sino el catálogo: no contiene suficientes
continuaciones factibles. Esta sección elimina el catálogo. La estrategia pasa a
ser la trayectoria misma, y la negociación ocurre mientras la planta avanza.

\subsection{Valor de misión y continuaciones ejecutables}

\begin{definicion}[Continuaciones ejecutables y valor de misión]
\label{def:continuaciones}
Sea $d$ un compromiso discreto (coalición, modo, contactos, ruta--tiempo,
reservas) y $Y$ el estado ampliado: estado físico y su estimación, estado de la
red, información disponible, continuación incumbente ya certificada, sus
márgenes y los recursos aún utilizables. El conjunto de continuaciones
ejecutables en horizonte $H$ es
\begin{equation}
  \label{eq:T-continuaciones}
  \Tcal(d,Y)=\bigl\{\tau=(x_{0:H},u_{0:H-1}):\ x_{k+1}=F(x_k,u_k,w_k),\
  (x_k,u_k,\lambda_k)\in\mathcal K_{\mathrm{phys}}(d),\ h_j(x_k)\ge0,\
  E_i(k)\ge E_i^{\min},\ u_i(k)\ \text{medible respecto de}\ \mathcal I_i(k),\
  x_H\in\mathcal G\bigr\},
\end{equation}
y, con $\mathcal H(\tau)$ el coste acumulado de la continuación, el valor del
compromiso es
\begin{equation}
  \label{eq:valor-mision}
  V_H^\star(d,Y)=\inf_{\tau\in\Tcal(d,Y)}\mathcal H(\tau),\qquad
  V_H^\star(d,Y)=+\infty\ \text{si}\ \Tcal(d,Y)=\emptyset .
\end{equation}
\end{definicion}

\begin{observacion}[Principio de coherencia de valor]
\label{obs:coherencia-valor}
Dos decisiones heterogéneas —reclutar, cambiar de contacto, desviar una ruta,
reemplazar un miembro— solo son comparables si se expresan como diferencias,
cotas o gradientes de la \emph{misma} función de valor
\eqref{eq:valor-mision}. Un coste de reclutamiento, una penalización de
\emph{wrench} o un precio de barrera no son por sí solos el valor de la misión.
La consecuencia no es retórica. Sobre dos mil mundos con veinte candidatos cada
uno, un proxy escalar de coste seleccionó un candidato \emph{inviable} en el
$76{,}15\%$ de los casos; el valor unificado, que devuelve $+\infty$ cuando
$\Tcal(d,Y)$ es vacío, en el $0\%$ \cite{mayorga2026compendio}. La infinitud de
\eqref{eq:valor-mision} es el mecanismo que impide que una relajación económica
invente una fuerza, un contacto o una ruta que la planta no puede realizar.
\end{observacion}

\subsection{Regla de \emph{commit} certificado y presupuesto de progreso}

El sistema distribuido rara vez dispone de $V_H^\star$ exacto. Dispone de
cotas.

\begin{proposicion}[\emph{Commit} certificado por cotas]
\label{pr:commit}
Sean $L_d(Y)\le V_H^\star(d,Y)\le U_d(Y)$ cotas locales, y sea $e:d\to d'$ una
revisión discreta con operador de reinicio $R_e$ y coste de transición
$C_e(Y)\le C_e^{+}(Y)$. Si
\begin{equation}
  \label{eq:commit}
  C_e^{+}(Y)+U_{d'}(R_eY)\;\le\;L_d(Y)-\delta,\qquad \delta>0,
\end{equation}
entonces la revisión mejora el valor real al menos en $\delta$:
$C_e(Y)+V_H^\star(d',R_eY)\le V_H^\star(d,Y)-\delta$.
\end{proposicion}

\begin{proof}
$C_e(Y)+V_H^\star(d',R_eY)\le C_e^{+}(Y)+U_{d'}(R_eY)\le L_d(Y)-\delta\le
V_H^\star(d,Y)-\delta$.
\end{proof}

\begin{observacion}[Conservadora por construcción]
\label{obs:commit-num}
La regla \eqref{eq:commit} no exige conocer el valor exacto y no puede producir
un \emph{commit} falso: en el banco reducido, $1\,383$ \emph{commits}
certificados y cero falsos \cite{mayorga2026compendio}; en veinte mil
instancias aleatorias con cotas de holgura uniforme en $[0,0{,}5]$, $4\,817$
\emph{commits} y cero falsos. El precio de esa garantía es que rechaza
revisiones buenas cuya mejora es menor que la holgura de las cotas: es la
misma asimetría que la condición de \emph{recourse} de \cref{pr:recourse}.
\end{observacion}

\begin{proposicion}[Presupuesto de progreso del testigo]
\label{pr:presupuesto-progreso}
Sea $W_k\ge0$ un presupuesto de valor asociado a la continuación incumbente
(el testigo de \cref{th:testigo}). Si cada intervalo físico de duración
$\Delta t_k$ consume al menos $\rho_0\Delta t_k$ y cada revisión aceptada
mejora al menos $\delta$, entonces
\begin{equation}
  \label{eq:presupuesto-progreso}
  0\le W_K\le W_0-\rho_0\sum_{k=0}^{K-1}\Delta t_k-N_{\mathrm{rev}}\,\delta ,
\end{equation}
luego el tiempo total está acotado por $W_0/\rho_0$ y el número de revisiones
por $W_0/\delta$. Con perturbación exógena acumulada $R$, la cota de tiempo se
desplaza a $(W_0+R)/\rho_0$.
\end{proposicion}

\begin{proof}
Ambos términos restan de $W$ y $W$ no puede ser negativo; despejar
\eqref{eq:presupuesto-progreso} con $N_{\mathrm{rev}}\ge0$ y con
$\sum\Delta t_k\ge0$ da cada cota. La perturbación suma a $W_0$.
\end{proof}

\begin{observacion}[Planear y mover dejan de alternar]
\label{obs:ejecutar-calculando}
Mientras el optimizador mejora una propuesta, la planta ejecuta la continuación
incumbente. \Cref{pr:presupuesto-progreso} es lo que hace legítima esa
simultaneidad: el progreso lo garantiza el testigo, no la convergencia
instantánea del solver, que es la lectura que ya daba \cref{pr:reloj-seguro}
para el reloj seguro. Lo que añade es que la cota es en \emph{valor de misión},
de modo que revisiones y tiempo físico se descuentan de la misma cantidad.
\end{observacion}

\subsection{Margen vectorial y \emph{recourse}}

No es correcto sumar márgenes con unidades físicas distintas.

\begin{definicion}[Margen vectorial de ejecutabilidad]
\label{def:margen-vectorial}
$m(d,Y)=(m_{\mathrm{sop}},m_{\mathrm{kin}},m_{\mathrm{wrench}},
m_{\mathrm{cont}},m_{\mathrm{seg}},m_{E},m_{\mathrm{ruta}},m_{\mathrm{info}})\ge0$:
soporte, compatibilidad cinemática (\cref{pr:ruta-ruedas}), reserva de
\emph{wrench}, contacto, seguridad, energía, ruta e información. Un
certificado está vigente mientras todas las componentes sean no negativas.
\end{definicion}

\begin{proposicion}[Condición suficiente de \emph{recourse}]
\label{pr:recourse}
Si cada componente decae a lo sumo linealmente, $\dot m_j\ge-\rho_j$, el tiempo
mínimo antes de perder algún certificado es
\begin{equation}
  \label{eq:T-hold}
  T_{\mathrm{hold}}=\min_{j:\rho_j>0}\frac{m_j(0)}{\rho_j},
\end{equation}
y con $T_{\mathrm{rec}}=T_{\mathrm{det}}+T_{\mathrm{com}}+T_{\mathrm{sol}}
+T_{\mathrm{apr}}+T_{\mathrm{adq}}+T_{\mathrm{commit}}$ el tiempo de
recuperación (detección, comunicación, resolución, aproximación, adquisición
de contacto y \emph{commit}), la condición
\begin{equation}
  \label{eq:recourse}
  T_{\mathrm{rec}}<T_{\mathrm{hold}}
\end{equation}
garantiza que los certificados siguen vigentes cuando el reemplazo queda
operativo.
\end{proposicion}

\begin{proof}
$m_j(t)\ge m_j(0)-\rho_jt>0$ para $t<m_j(0)/\rho_j$; el mínimo sobre $j$ da
\eqref{eq:T-hold}, y \eqref{eq:recourse} coloca el fin del reemplazo antes.
\end{proof}

\begin{observacion}[Qué cuantifica y qué no]
\label{obs:recourse-num}
La condición no afirma que toda falla sea recuperable: mide cuándo la
ejecución incumbente compra tiempo suficiente para una revisión. En cinco mil
fallos simulados, ningún caso certificado resultó inseguro; se certificó el
$11{,}64\%$ mientras el $27{,}72\%$ era realmente recuperable
\cite{mayorga2026compendio}. La condición es conservadora por un factor
próximo a $2{,}4$, que es el precio de una cota lineal sobre cada margen. La
misma cantidad $T_{\mathrm{hold}}$ es la que \cref{th:prima} necesita del lado
de la reserva: la prima paga tiempo de recuperación, y \eqref{eq:recourse}
dice cuánto tiempo hay que comprar.
\end{observacion}

\subsection{La trayectoria como estrategia}

Un AMR libre es un jugador. Una coalición en \textsf{cargo} ya cerrada es
externamente un solo jugador compuesto, porque sus apoyos no pueden separarse
durante el transporte y su radio de giro sale de \cref{pr:ruta-ruedas}. Para
cada jugador $a$, la estrategia es la continuación
\begin{equation}
  \label{eq:tau-jugador}
  \tau_a(\cdot)=\bigl(q_a(\cdot),u_a(\cdot)\bigr),\qquad \cdot\in[t,t+H],
\end{equation}
sin \emph{waypoint} impuesto: la única propuesta inicial es la recta hacia el
objetivo o la continuación desplazada del ciclo anterior. Cada jugador publica
una ocupación prevista $\rho_a(x,\tau)\ge0$ que incorpora su huella —para
\textsf{cargo}, la carga y el conjunto rígido de soportes—, la ocupación
agregada es $\rho=\sum_a\rho_a$, y un precio dual local de congestión se
actualiza en tiempo de cálculo $\sigma$ como
\begin{equation}
  \label{eq:precio-congestion}
  \frac{\mathrm d\mu}{\mathrm d\sigma}(x,\tau)=\Pi_{\mathbb R_+}\bigl(\rho(x,\tau)-\rho^{\max}(x,\tau)\bigr),
\end{equation}
que no asigna turnos: altera gradientes y deja que espera, desvío, aceleración
o giro emerjan como respuesta. El coste individual usado en los ensayos es
\begin{equation}
  \label{eq:coste-trayectoria}
  J_a(\tau_a,\tau_{-a})=w_TT_a+w_f\lVert q_a(t+H)-q_a^{d}\rVert^2
  +\int_t^{t+H}\bigl(w_u\lVert u_a\rVert^2+w_\kappa\kappa_a^2+B_{\mathrm{obs}}
  +C_{\mathrm{cong}}+C_{\mathrm{int}}+C_{\mathrm{phys}}\bigr)\,\mathrm d\tau ,
\end{equation}
donde el término temporal es esencial: sin él, detenerse indefinidamente puede
resultar barato.

Para pocas coaliciones grandes no es correcto usar la derivada infinitesimal de
la congestión. Sea $\mathcal H(\rho)=\int_t^{t+H}\!\!\int_\Omega
H(\rho(x,\tau))\,\mathrm dx\,\mathrm d\tau$.

\begin{teorema}[Externalidad atómica exacta]
\label{th:externalidad-atomica}
Si el coste de congestión de cada jugador es su contribución marginal exacta,
$C_a^{\mathrm{cong}}=\mathcal H(\rho_{-a}+\rho_a)-\mathcal H(\rho_{-a})$, y las
interacciones restantes son simétricas por pares, $\Phi_{ab}=\Phi_{ba}$,
entonces
\begin{equation}
  \label{eq:potencial-trayectorias}
  P(\tau)=\sum_a J_a^{0}(\tau_a)+\mathcal H\Bigl(\sum_a\rho_a\Bigr)+\sum_{a<b}\Phi_{ab}(\tau_a,\tau_b)
\end{equation}
es un potencial exacto del juego de trayectorias: para toda desviación
unilateral $\tau_a\to\tau_a'$, $J_a(\tau_a',\tau_{-a})-J_a(\tau_a,\tau_{-a})
=P(\tau_a',\tau_{-a})-P(\tau_a,\tau_{-a})$, con igualdad de diferencias
finitas y no solo de gradientes.
\end{teorema}

\begin{proof}
La diferencia de $J_a$ es $\Delta J_a^{0}+\mathcal H(\rho_{-a}+\rho_a')
-\mathcal H(\rho_{-a}+\rho_a)+\sum_{b\ne a}\Delta\Phi_{ab}$, porque el término
$\mathcal H(\rho_{-a})$ se cancela. La diferencia de $P$ es exactamente la
misma suma: los $J_b^{0}$ y los $\Phi_{bc}$ con $b,c\ne a$ no cambian.
\end{proof}

\begin{observacion}[Verificación y lectura]
\label{obs:externalidad-num}
Sobre trescientas desviaciones unilaterales aleatorias con $H$ convexa no
cuadrática y ocupaciones gaussianas, la discrepancia máxima entre
$\Delta J_a$ y $\Delta P$ es $1{,}2\times10^{-14}$. El teorema hace del
planificador de trayectorias una \emph{rama} del megajuego de
\cref{th:potencial}: dentro de una rama discreta fija, el equilibrio de
\eqref{eq:potencial-trayectorias} sobre un conjunto convexo coincide con el
óptimo social de esa rama. La equivalencia no se extiende por sí sola a
obstáculos no convexos, cambios de contacto ni combinatoria discreta, que es
exactamente lo que \cref{obs:homotopia} recoge.
\end{observacion}

Numéricamente la trayectoria se representa por $K$ coeficientes B-spline,
$q_a(\tau)=\sum_{k=1}^{K}B_k(\tau)c_{a,k}$, o por una secuencia de controles
$(v_k,\omega_k)$; en tiempo de cálculo los coeficientes se deforman por
$\dot c_a=-\Gamma_aK_a\nabla_{c_a}P$, y el campo deseado
$v_a^{\mathrm{game}}=-M_a^{-1}\nabla_xC_a$ se ejecuta dentro de la envolvente
física $u_a=U_a(q)(v_a^{\mathrm{game}})$. La versión limpia no corrige a
posteriori: ruedas, curvatura y huella entran en $K_a$, de modo que la
diferencia entre lo pedido y lo ejecutado es pequeña.

\begin{proposicion}[Dinámica Smith en caja]
\label{pr:smith-caja}
Para una acción escalar $a_i\in[-T_i,T_i]$ y gradiente $\gamma_i=\partial
J/\partial a_i$, la realización
\begin{equation}
  \label{eq:smith-caja}
  \dot a_i=2\lambda_iT_i\Bigl[(T_i-a_i)[-\gamma_i]_+-(T_i+a_i)[\gamma_i]_+\Bigr]
\end{equation}
deja la caja invariante y desciende $J$: en $a_i=T_i$ desaparece el término
que empuja hacia arriba y en $a_i=-T_i$ el que empuja hacia abajo.
\end{proposicion}

\begin{proof}
En $a_i=T_i$, $\dot a_i=-2\lambda_iT_i\cdot2T_i[\gamma_i]_+\le0$; en
$a_i=-T_i$, $\dot a_i=2\lambda_iT_i\cdot2T_i[-\gamma_i]_+\ge0$. Para el
descenso, si $\gamma_i\ge0$ entonces
$\gamma_i\dot a_i=-2\lambda_iT_i(T_i+a_i)\gamma_i^2\le0$, y si $\gamma_i<0$
entonces $\gamma_i\dot a_i=-2\lambda_iT_i(T_i-a_i)\gamma_i^2\le0$; en veinte
mil pasos con gradiente aleatorio la acción no sale de la caja.
\end{proof}

\begin{observacion}[Extragradiente: por qué]
\label{obs:extragradiente-silla}
Con componente rotacional, Euler simultáneo diverge. En la silla bilineal
$F(z)=(y,-x)$ el paso explícito multiplica la norma por $(1+h^2)^{1/2}$ en
cada iteración y el extragradiente por $(1-h^2+h^4)^{1/2}$: con $h=0{,}5$ y
$69$ pasos, $\times2\,205$ frente a $\times7{,}7\times10^{-4}$, que es el
régimen del contraste reportado en \cite{mayorga2026compendio}
($\times2\,180$ frente a $1{,}43\times10^{-3}$). Es el mismo argumento de
\cref{sec:comunicacion}: valida el mecanismo numérico, no una garantía global
del megajuego no lineal.
\end{observacion}

\begin{proposicion}[Intención publicada con incertidumbre]
\label{pr:convolucion-intencion}
Si el jugador $b$ publica su pose futura como $Q_b(\tau)\sim\mathcal N(\hat
q_b(\tau),\Sigma_b(\tau))$ y la huella se representa con núcleo gaussiano
$K_S(d)=\exp(-\tfrac12d^{\!\top}S^{-1}d)$, la ocupación esperada es exacta:
\begin{equation}
  \label{eq:convolucion}
  \mathbb E\bigl[K_S(q-Q_b)\bigr]=\frac{|S|^{1/2}}{|S+\Sigma_b|^{1/2}}
  \exp\Bigl(-\tfrac12(q-\hat q_b)^{\!\top}(S+\Sigma_b)^{-1}(q-\hat q_b)\Bigr).
\end{equation}
\end{proposicion}

\begin{proof}
Es la convolución de dos gaussianas; con $400\,000$ muestras Monte Carlo la
forma cerrada y la media muestral coinciden en $0{,}2746$ frente a $0{,}2748$.
\end{proof}

\begin{observacion}[La incertidumbre entra en el precio, no en el margen]
\label{obs:convolucion-uso}
\Cref{pr:convolucion-intencion} es lo que permite que la covarianza de
\cref{co:tensado} entre en el campo de ocupación en lugar de en un ensanche
fijo de la huella: un vecino incierto ocupa más espacio pero con menos
densidad, y el precio \eqref{eq:precio-congestion} lo cobra en proporción. En
el banco reducido, incorporar la incertidumbre futura al mercado bajó la
colisión Monte Carlo a $\Delta t=1{,}2$~s del $13{,}11\%$ al $4{,}06\%$ por un
$2{,}8\%$ más de recorrido \cite{mayorga2026compendio}.
\end{observacion}

\begin{proposicion}[Conectividad acumulada antes del conflicto]
\label{pr:conectividad-acumulada}
Sea $e_I$ el desacuerdo de una señal distribuida de intención o precio con
$\dot e_I=-cL(t)e_I$ y $e_I\perp\mathbf 1$. Entonces
$\lVert e_I(t)\rVert\le\exp\bigl(-c\int_{t_0}^{t}\lambda_2(L(s))\,\mathrm ds\bigr)\lVert e_I(t_0)\rVert$,
y si el juego necesita un tiempo $T_G$ tras alcanzar tolerancia $\varepsilon_I$
antes del instante de conflicto $t_c$, basta
\begin{equation}
  \label{eq:conectividad-acumulada}
  \int_{t_0}^{t_c-T_G}\lambda_2(L(s))\,\mathrm ds\;\ge\;\frac1c\log\frac{\lVert e_I(t_0)\rVert}{\varepsilon_I}.
\end{equation}
\end{proposicion}

\begin{proof}
$\tfrac{\mathrm d}{\mathrm dt}\lVert e_I\rVert^2=-2c\,e_I^{\!\top}L e_I\le-2c\lambda_2\lVert e_I\rVert^2$
por Rayleigh en $\mathbf 1^\perp$; integrar y exigir
$\lVert e_I(t_c-T_G)\rVert\le\varepsilon_I$.
\end{proof}

\begin{observacion}[Lo que se defiende]
\label{obs:conectividad-mecanismo}
No hace falta afirmar que toda congestión aumenta $\lambda_2$. El mecanismo
defendible es que la aproximación a un cuello crea enlaces de proximidad, y
si la conectividad acumulada supera el presupuesto
\eqref{eq:conectividad-acumulada}, información y precios convergen antes de la
interacción física. En $512$ ensayos la condición no dio ningún falso positivo,
siendo conservadora \cite{mayorga2026compendio}. Es la misma $\lambda_2$ que
gobierna la percepción en \cref{sec:estimacion}: un solo número controla
cuándo el mapa y cuándo el precio están listos.
\end{observacion}

\subsection{Dos relojes: separación de escalas}

En cada instante físico $t_k$ se desplaza la solución anterior como arranque,
se reciben las últimas propuestas, se ejecutan unas pocas iteraciones del
juego en tiempo de cálculo, se aplica el primer control de la continuación
vigente y la planta avanza. Abstractamente,
\begin{equation}
  \label{eq:perturbacion-singular}
  \dot X=f\bigl(X,\pi(z,X)\bigr),\qquad \varepsilon\dot z=-F(z;X),\qquad
  \varepsilon=\frac{T_{\mathrm{game}}}{T_{\mathrm{phys}}}\ll1 .
\end{equation}

\begin{observacion}[Seguimiento cuasiestático y sus cuatro condiciones]
\label{obs:escalas}
Si para $X$ congelado $F(z;X)=0$ tiene solución única $z^\star(X)$
uniformemente exponencialmente estable y $z^\star$ es Lipschitz, la teoría de
perturbación singular da cotas del tipo
$\lVert z(t)-z^\star(X(t))\rVert\le c_1e^{-t/\varepsilon}\lVert z(0)-z^\star(X(0))\rVert+c_2\varepsilon$,
y con errores de comunicación, solver, modelo y muestreo,
$\lVert z-z^\star\rVert\le C(\varepsilon+e_{\mathrm{com}}+e_{\mathrm{sol}}+e_{\mathrm{mod}}+e_{\mathrm{mue}})$.
Que el juego sea mucho más rápido que la planta no implica optimalidad: permite
seguir con error pequeño el equilibrio de la rama que el solver resuelve, y si
esa rama tiene el mínimo local equivocado, acelerar converge más rápido al
mínimo equivocado. Para afirmar $J\approx J^\star_{\mathrm{Mega}}$ hacen
falta además: equilibrio rápido equivalente al óptimo de la rama
(\cref{th:externalidad-atomica} con convexidad), cobertura de las ramas
relevantes, modelo predictivo coherente con la planta, y horizonte completo o
valor terminal de Bellman con error acotado. Un solver infinitamente rápido
con terminal miope resuelve un problema miope.
\end{observacion}

\begin{observacion}[Evidencia de mecanismo]
\label{obs:escalas-num}
En diez mundos pareados alrededor de un cruce, cambiar solo el número de
iteraciones internas de negociación pasó el éxito seguro de $3/10$ a $9/10$:
seis mundos mejoraron y ninguno empeoró, y el McNemar exacto con seis pares
discordantes en la misma dirección es $p=2\cdot2^{-6}=0{,}03125$
\cite{mayorga2026compendio}. Es evidencia del mecanismo de
\eqref{eq:perturbacion-singular}, no una demostración del teorema; la planta
es 2D reducida y la energía es un proxy.
\end{observacion}

\begin{observacion}[Cuatro relojes, un solo juego]
\label{obs:cuatro-relojes}
La política recomendada es
$T_{\mathrm{wrench}}<T_{\mathrm{path}}<T_{\mathrm{phys}}<T_{\mathrm{coalición}}$:
el mercado de \emph{wrench} se limpia dentro de cada paso de trayectoria, la
trayectoria dentro de cada paso físico, y la coalición cambia en la escala más
lenta. No son cuatro juegos: son cuatro constantes de tiempo de la misma
estrategia híbrida, y su orden es la condición de \cref{obs:escalas} aplicada
por capas.
\end{observacion}

\subsection{Qué significa «la mejor ruta»}

\begin{definicion}[Jerarquía de referencias]
\label{def:jerarquia-optimos}
El óptimo completo $J^\star_{\mathrm{Mega}}=\min_{d,X,U,\lambda}J$ sujeto a
reclutamiento, modos, contactos, dinámica, ruedas, \emph{wrench}, fricción,
seguridad, energía, comunicación y entrega es híbrido y no convexo. Con
coaliciones y modos congelados, una búsqueda A* conjunta espacio--tiempo sobre
un lattice finito $\mathcal L$ con heurística admisible devuelve
$J^\star_{\mathcal L}=\min_{\tau\in\mathcal L}J(\tau)$, y como $\mathcal L$
restringe el continuo, $J^\star_{\mathrm{Mega}}\le J^\star_{\mathcal L}$. Un
optimizador central continuo con toda la información devuelve el mejor
programa no lineal encontrado, no un óptimo certificado. La brecha
$\mathrm{Gap}=(J_{\mathrm{game}}-J_{\mathrm{oracle}})/J_{\mathrm{oracle}}$
debe declarar siempre cuál de los tres es $J_{\mathrm{oracle}}$.
\end{definicion}

\begin{observacion}[Lectura del banco de treinta mundos]
\label{obs:lectura-r10}
El banco de \cref{obs:ensayo-r10} da la lectura: en la capa continua, el
megajuego alcanza calidad numérica indistinguible de un MPC distribuido por
intercambio de trayectorias —diferencias medianas de tiempo, distancia y
energía de $0{,}000\%$, $+0{,}028\%$ y $+0{,}032\%$ en veintisiete mundos
comunes, ninguna significativa tras Holm— y queda a $+1{,}6\%$ de tiempo del
mejor central encontrado. La afirmación correcta no es «un MPC distribuido
mejor» sino que el mismo mecanismo de precios, estrategias y
\emph{commit} sigue siendo válido cuando la decisión pasa de trayectoria a
coalición, contacto, \emph{wrench}, reserva o reemplazo, que el MPC
distribuido no modela. Y en el pasillo ancho, donde la geometría no exige
negociación, un método local ligero es más rápido: no hay superioridad
universal que reclamar.
\end{observacion}

\begin{observacion}[La homotopía es parte del compromiso]
\label{obs:homotopia}
El escenario radial de \cref{obs:ensayo-r10} muestra la debilidad de la
deformación continua: puede quedarse atrapada en una familia topológica que no
lleva al mejor camino. La respuesta es elevar la clase de homotopía $h$ a
componente del compromiso discreto, $d=(\Ccal_R,\Ccal_A,\mu,\sigma,h,b,\xi)$:
dentro de una clase, \eqref{eq:potencial-trayectorias} deforma $\tau$; entre
clases, el grafo de revisión de \cref{obs:cuatro-grafos} admite una arista
$h\to h'$ solo si existe una transición físicamente ejecutable con testigo de
seguridad. Optimización rápida dentro de la rama y salto discreto entre ramas:
la pregunta teórica pasa de «cuántas iteraciones de gradiente» a «qué soporte
de revisiones conecta todas las ramas relevantes con coste controlado», que es
la pregunta de \cref{th:saltos}.
\end{observacion}

\begin{observacion}[Reserva robusta, no más peso]
\label{obs:reserva-robusta}
Los ensayos continuos dan separaciones mínimas de $3{,}1$~cm
(\cref{obs:ensayo-r6}), demasiado pequeñas para un margen industrial. La
corrección no es subir el peso de la penalización sino imponer una reserva:
para dos huellas, $h_{ab}=d_{ab}(X)-d^{\mathrm{rob}}_{\mathrm{safe}}\ge0$ con
$d^{\mathrm{rob}}_{\mathrm{safe}}=d_{\mathrm{safe}}+\overline{\Delta d}+\delta_{\mathrm{mue}}$,
que es \cref{co:tensado} con el término de muestreo, y si el par tiene grado
relativo $r>1$ respecto de $h$, la restricción se rederiva como barrera de
orden superior en lugar de reutilizar una barrera de aceleración como si el
par entrase en primer orden.
\end{observacion}

\subsection{Reparar localmente, ampliar cuando no basta, y elegir el solver dentro de la rama}

Tras una perturbación que afecta a una parte limitada de la flota hay tres
decisiones distintas: qué región renegocia ($S$), cuántos agentes cambian
juntos ($h$) y qué precisión basta ($\varepsilon$). La batería R11
\cite{mrob2026r11} las separa con problemas reducidos, referencias
calculables y controles negativos; lo que sigue es lo que sobrevive.

\begin{proposicion}[Certificado de calidad por residual]
\label{pr:certificado-gradiente}
Si $\Phi$ es diferenciable y $m$-fuertemente convexa sin restricciones,
$0\le\Phi(x)-\Phi(x^\star)\le\lVert\nabla\Phi(x)\rVert^2/(2m)$. Si $\Phi$ es
convexa sobre una caja $K$ y $U\in K$, la brecha de Frank--Wolfe
$G_{\mathrm{FW}}(U)=\max_{Z\in K}\nabla\Phi(U)^{\!\top}(U-Z)$, que en una caja
se calcula componente a componente, cumple
$\Phi(U)-\Phi(U^\star)\le G_{\mathrm{FW}}(U)$.
\end{proposicion}

\begin{proof}
Convexidad fuerte: $\Phi(y)\ge\Phi(x)+g^{\!\top}(y-x)+\tfrac m2\lVert y-x\rVert^2$;
el mínimo en $y$ del lado derecho es $\Phi(x)-\lVert g\rVert^2/(2m)$, y se
evalúa en $y=x^\star$. Con caja: $\Phi(U^\star)\ge\Phi(U)+g^{\!\top}(U^\star-U)$
y maximizar $g^{\!\top}(U-Z)$ sobre $K$ da la cota.
\end{proof}

\begin{observacion}[Reparación local certificada: cuándo ahorra y cuándo no]
\label{obs:reparacion-local}
Sobre un potencial cuadrático en malla de cuatro vecinos,
$\Phi=\tfrac m2\lVert x\rVert^2+\tfrac12\sum_{ij\in E}(x_i-x_j)^2-b^{\!\top}x$,
una perturbación local de la demanda se repara actualizando solo una región que
crece hasta que \cref{pr:certificado-gradiente} cae por debajo de $10^{-6}$,
frente a actualizar toda la red con el mismo operador. Con $m=1$ y doscientas
cincuenta perturbaciones, todas certificadas: en $N=100$ la región cuesta un
$4{,}95\%$ más de actualizaciones; en $N=1\,024$, un $83\%$ menos; en
$N=10\,000$, un $98{,}1\%$ menos, con intervalo bootstrap $97{,}8$--$98{,}2\%$
y brecha real máxima $3{,}8\times10^{-7}$ \cite{mrob2026r11}. Una réplica
propia con protocolo de expansión distinto da $65\%$ y $94{,}6\%$ menos en los
mismos tamaños. Los controles negativos son la parte que importa: con
perturbación difusa la región activa toda la red y gasta un $18\%$ más; con
$m=0{,}01$ ninguna de cincuenta perturbaciones se certifica dentro de radio
doce (réplica propia: cero de diez, brecha $2{,}7$). La hipótesis defendible
es «localidad útil cuando perturbación y acoplamiento permiten una reparación
local», que es la caída exponencial de la sensibilidad en programas
estructurados por grafos \cite{shin2022exponential}, no «siempre basta un
vecindario».

Dos precisiones que evitan sobreclaims. La mejor respuesta de cada nodo,
$x_i^{+}=(b_i+\sum_{j\in N_i}x_j)/(m+\deg i)$, es exactamente un paso de
gradiente precondicionado por la diagonal —discrepancia $5{,}5\times10^{-16}$—,
de modo que la ganancia viene de localizar la reparación y no del nombre del
algoritmo. Y ahorrar actualizaciones no es ahorrar tiempo: en $N=10\,000$ la
reparación local tardó $5{,}33$~ms frente a $5{,}06$~ms del método global y
$0{,}96$~ms de una factorización ya preparada. Actualizaciones, mensajes, CPU y
latencia física son cuatro magnitudes; \cref{obs:coherencia-valor} exige
declarar cuál se reduce.
\end{observacion}

\begin{observacion}[Más cálculo reduce el retraso; no corrige el soporte]
\label{obs:mas-calculo}
Sobre cincuenta óptimos móviles de una red de $256$ variables con un salto no
anunciado, pasar de una a dieciséis iteraciones por muestra reduce el error
cuadrático medio respecto del óptimo instantáneo un $99{,}49\%$ en mediana
pareada, mejorando las cincuenta realizaciones; pero ciento veintiocho
iteraciones con la región congelada a radio dos quedan peor que una sola
iteración sobre toda la red ($4{,}9\times10^{-2}$ frente a
$2{,}6\times10^{-2}$) \cite{mrob2026r11}. Es la segunda mitad de
\cref{obs:escalas} medida: para un operador contractivo de factor $\rho$,
$\lVert e_{k+1}\rVert\le\rho^{K}(\lVert e_k\rVert+\lVert x^\star_{k+1}-x^\star_k\rVert)$
explica el retraso que $K$ elimina, y una región truncada añade un sesgo que
$K$ no toca. Lo que hay que medir en la planta es rapidez \emph{y} suficiencia
del dominio de negociación.
\end{observacion}

\begin{observacion}[El solver dentro de la rama: un resultado que cambia el diseño]
\label{obs:solver-en-rama}
Sobre cincuenta programas predictivos de una carga con cuatro AMR en rama
longitudinal fija —doce pasos de $0{,}2$~s, cuarenta y ocho variables, límites
heterogéneos de actuación, un AMR con autoridad reducida, mismo coste, mismo
estado y misma tolerancia $G_{\mathrm{FW}}\le10^{-7}$—, un programa cuadrático
central, mejores respuestas cíclicas del juego potencial y un ADMM de consenso
de dos bloques \cite{boyd2011admm} llegan al mismo óptimo (diferencia máxima
$4\times10^{-13}$), pero con costes muy distintos: mejores respuestas,
trescientas rondas y $630$~ms medianos; ADMM, treinta y dos rondas y
$1{,}87$~ms; el central, $3{,}0$~ms \cite{mrob2026r11}. Una réplica propia
reproduce la parte robusta del resultado y matiza la otra: el mismo óptimo a
$10^{-14}$, un factor $137$ de cómputo a favor de ADMM ($10{,}6$ frente a
$1\,450$~ms), pero rondas comparables ($116$ frente a $98$): el cociente de
rondas depende del parámetro $\rho$ y de la implementación, y la diferencia de
cómputo no, porque la mejor respuesta resuelve un programa interno en cada
bloque y el paso de ADMM es proximal cerrado,
$U_i^{+}=\Pi_{[-\bar u_i,\bar u_i]}\bigl(\rho(S_i-Y_i)/(r_i+\rho)\bigr)$.

La lectura correcta es esta. El juego es la estructura de decisiones, pagos y
precios que hace comparables coalición, contacto, ruta y reemplazo bajo la
misma función de valor; dentro de una rama convexa, cualquier solver que
alcance el mismo equilibrio variacional vale, y debe elegirse el eficiente.
No hay obligación científica de que la capa continua use mejores respuestas o
Smith para «parecer más de juegos»: los juegos poblacionales ya se combinan con
control predictivo distribuido en la literatura \cite{barreiro2019evolutionary},
y \cref{th:externalidad-atomica} garantiza que el equilibrio de la rama es el
mismo sea cual sea el método. Es también la respuesta a
\cref{obs:lectura-r10}: la paridad numérica con el MPC distribuido no es una
coincidencia, es que resuelven el mismo problema, y el aporte está en lo que
el MPC distribuido no modela.
\end{observacion}

\begin{observacion}[Lo que R11 no demuestra]
\label{obs:r11-alcance}
Mil quinientas cincuenta realizaciones de problema, no misiones: juegos
cuadráticos en red, asignación atómica de ocho AMR, óptimos móviles, frenados
longitudinales y programas predictivos lineales. Ninguno integra los
mecanismos en una misma simulación física de flota abierta, ni compara con
un MPC distribuido equipado con la misma reparación adaptativa. La hipótesis
de sistema que queda delimitada, y que se falsaría con el mismo modelo,
información y presupuesto, es: calidad de misión comparable a un MPC
distribuido competente, con menos trabajo y disrupción de adaptación ante
eventos localmente reparables.
\end{observacion}

\subsection{Tres cotas del ciclo de vida}

\begin{lema}[Error de valoración de propuestas]
\label{le:valoracion-2c}
Si $|\hat\Psi(B)-\Psi(B)|\le c$ para toda propuesta $B$ y
$\hat B\in\arg\min\hat\Psi$, entonces $\Psi(\hat B)\le\Psi(B^\star)+2c$.
\end{lema}

\begin{proof}
$\Psi(\hat B)\le\hat\Psi(\hat B)+c\le\hat\Psi(B^\star)+c\le\Psi(B^\star)+2c$.
En veinte mil sorteos con veinte propuestas y $c=0{,}3$ el exceso máximo
observado es $0{,}546<0{,}6$.
\end{proof}

\begin{observacion}[Determinista antes que Gibbs]
\label{obs:valoracion-gibbs}
\Cref{le:valoracion-2c} es el término determinista que complementa a
\cref{th:gibbs}: un solver interno de precisión $c$ entrega, antes de
cualquier aleatorización, un ganador a lo sumo $2c$ peor en valor real. La
condición $\beta>2c$ sobre la temperatura inversa hace que la concentración
de Gibbs no deshaga lo que la precisión del solver garantiza.
\end{observacion}

\begin{observacion}[Salto de precio en lugar de reinicio]
\label{obs:salto-precio}
Cuando un miembro entra o sale, el mercado de \emph{wrench} no se reinicia: se
aplica un paso de Newton sobre el mapa agregado,
$p^{+}=p^{-}+(J_W^{+})^{-1}\bigl[w_d^{+}-W^{+}(p^{-})\bigr]$, y se continúa. En
el banco reducido, las medianas de iteraciones hasta limpiar fueron $156{,}5$
en frío, $143$ con arranque caliente y $80$ con el salto
\cite{mayorga2026compendio}. Es la versión de evento de \cref{th:saltos}: la
red de precios de \cref{sec:nucleo} tiene forma cerrada en el caso lineal, y
el salto es exacto ahí y aproximado en el no lineal.
\end{observacion}

\begin{observacion}[Propuestas conjuntas frente a ranking individual]
\label{obs:bundles}
La unidad discreta de negociación no debe ser un AMR aislado sino la propuesta
conjunta $(\Ccal,m,r,h)$: miembros, modo físico, geometría de contacto y clase
de ruta. Un AMR cercano puede perder frente a uno lejano si su incorporación
empeora la geometría de \emph{wrench}, obliga a un corredor congestionado o
reduce la reserva. En el banco reducido, un ranking escalar por AMR alcanzó el
$78{,}3\%$ de éxito frente al $100\%$ de las propuestas conjuntas, y cuando
funcionó costó $1{,}394$ veces más y usó $0{,}91$ AMR extra
\cite{mayorga2026compendio}. Es la no superaditividad de \cref{sec:standby}
vista desde el reclutamiento: la complementariedad no se captura sumando
valores individuales.
\end{observacion}
